\scanpage{74}
\mspage{\handwritten{41}}
après le changement de base $X'\longrightarrow X$, et on applique la partie déjà prouvée de 9.1. Le morphisme $u:F\longrightarrow G$ répond alors à la question, cqfd.
\margin{\handwritten{\tikz\draw (0,0) .. controls (-.12,.19) and (-.12,.54) .. (0,.86);}}

\uline{Théorème} 9.4. Soit $f:X'\longrightarrow X$ un morphisme surjectif de préschéma. On suppose que l'une des hypothèses suivantes est vérifiée:

a) $f$ est entier.

b) $f$ est propre.

c) $f$ est plat et localement de présentation finie.

d) $X$ est discret (p.ex. le spectre d'un corps).

% slanted
\margin{\handwritten{pas d'interligne suppl.}}
Alors $f$ est un \uline{morphisme de descente effective} pour la catégorie fibrée $\unsure{\handwritten{\mathfrak F}|\mathfrak F|\mathcal F}$ sur (Sch) des faisceaux étales sur des préschémas variables, i.e. le foncteur $f^{*}$ induit une \uline{équivalence} de la catégorie \margin{\handwritten{\tikz\draw (0,0) .. controls (.08,.23) and (.19,.32) .. (.33,.22);}}des faisceaux sur $X$ avec la catégorie des faisceaux sur $X'$, \uline{munis} d'une donnée de descente relativement à $f:X'\longrightarrow X$.

\margin{\handwritten{\tikz\draw (0,0) .. controls (-.18,.15) and (-.17,.43) .. (0,.63);}}
9. 4.1. Soit $F'$ un faisceau sur $X'$, muni d'une donnée de descente relativement à $f:X'\to X$. On définit alors un foncteur
\margin{\handwritten{\tikz[baseline=(w.base)]{\node (w) {centrer};\draw (w.center) ellipse (.75 and .27);}}}
\[
\deleted{\handwritten{\unsure{\longleftarrow}}}\quad\handwritten{\longleftarrow}\quad G:X_{\mathrm{ét}}^{\unsure{o|0|\circ}}\longrightarrow(\mathrm{Ens})
\]
\margin{\handwritten{\tikz[baseline=(w.base)]{\node[align=center] (w) {\unsure{replacer} à\\la ligne};\draw (w.center) ellipse (.87 and .53);}}}
\handwritten{\(\longleftarrow\)}
par la formule
\margin{\handwritten{\tikz[baseline=(w.base)]{\node (w) {centrer};\draw (w.center) ellipse (.75 and .27);}}}
\[
G(Y)=\operatorname{Ker}\left(F'(Y')\rightrightarrows F''(Y'')\right).
\]
On constate aussitôt que $G$ est un \uline{faisceau} pour la topologie étale, d'autre part on a un homomorphisme canonique injectif $G\to f_{*}(F')$,
% archive sheet number: 73
